\documentclass{beamer}

\usepackage[T2A]{fontenc}
\usepackage[utf8]{inputenc}
\usepackage[english,russian]{babel}
\pdfmapfile{+sansmathaccent.map}


\mode<presentation>
{
	\usetheme{Warsaw} % or try Darmstadt, Madrid, Warsaw, Rochester, CambridgeUS, ...
	\usecolortheme{seahorse} % or try seahorse, beaver, crane, wolverine, ...
	\usefonttheme{serif}  % or try serif, structurebold, ...
	\setbeamertemplate{navigation symbols}{}
	\setbeamertemplate{caption}[numbered]
} 


%%%%%%%%%%%%%%%%%%%%%%%%%%%%
% itemize settings


%%%%%%%%%%%%%%%%%%%%%%%%%%%%
% itemize settings

\definecolor{myhotpink}{RGB}{255, 80, 200}
\definecolor{mywarmpink}{RGB}{255, 60, 160}
\definecolor{mylightpink}{RGB}{255, 80, 200}
\definecolor{mypink}{RGB}{255, 30, 80}
\definecolor{mydarkpink}{RGB}{155, 25, 60}

\definecolor{mypaleblue}{RGB}{240, 240, 255}
\definecolor{mylightblue}{RGB}{120, 150, 255}
\definecolor{myblue}{RGB}{90, 90, 255}
\definecolor{mygblue}{RGB}{70, 110, 240}
\definecolor{mycourseblue}{RGB}{120, 140, 255}
\definecolor{mycourselightblue}{RGB}{245, 250, 255}
\definecolor{mydarkblue}{RGB}{0, 0, 180}
\definecolor{myblackblue}{RGB}{40, 40, 120}

\definecolor{mygreen}{RGB}{0, 200, 0}
\definecolor{mygreen2}{RGB}{245, 255, 230}
\definecolor{mydarkgreen}{RGB}{0, 120, 0}
\definecolor{mydarkgreen}{RGB}{0, 120, 0}


\definecolor{mygray}{gray}{0.8}
\definecolor{mydarkgray}{RGB}{80, 80, 160}
\definecolor{mylightgray}{RGB}{160, 160, 160}

\definecolor{mydarkred}{RGB}{160, 30, 30}
\definecolor{mylightred}{RGB}{255, 150, 150}
\definecolor{myred}{RGB}{200, 110, 110}
\definecolor{myblackred}{RGB}{120, 40, 40}

\definecolor{mypink}{RGB}{255, 30, 80}
\definecolor{myhotpink}{RGB}{255, 80, 200}
\definecolor{mywarmpink}{RGB}{255, 60, 160}
\definecolor{mylightpink}{RGB}{255, 80, 200}
\definecolor{mydarkpink}{RGB}{155, 25, 60}
\definecolor{mywhitepink}{RGB}{255, 240, 240}

\definecolor{mydarkcolor}{RGB}{60, 25, 155}
\definecolor{mylightcolor}{RGB}{130, 180, 250}



\setbeamertemplate{itemize items}[default]

\setbeamertemplate{itemize item}{\color{myblackblue}$\blacksquare$}
\setbeamertemplate{itemize subitem}{\color{mygblue}$\blacktriangleright$}
\setbeamertemplate{itemize subsubitem}{\color{mygray}$\blacksquare$}

\setbeamercolor{palette quaternary}{fg=white,bg=mycourseblue}
\setbeamercolor{titlelike}{parent=palette quaternary}

\setbeamercolor{palette quaternary2}{fg=black,bg=mycourselightblue}
\setbeamercolor{frametitle}{parent=palette quaternary2}

\setbeamerfont{frametitle}{size=\Large,series=\scshape}
\setbeamerfont{framesubtitle}{size=\normalsize,series=\upshape}





%%%%%%%%%%%%%%%%%%%%%%%%%%%%
% block settings

\setbeamercolor{block title}{bg=red!30,fg=black}

\setbeamercolor*{block title example}{bg=mygreen!40!white,fg=black}

\setbeamercolor*{block body example}{fg= black, bg= mygreen2}


%%%%%%%%%%%%%%%%%%%%%%%%%%%%
% URL settings
\hypersetup{
	colorlinks=true,
	linkcolor=blue,
	filecolor=blue,      
	urlcolor=blue,
}

%%%%%%%%%%%%%%%%%%%%%%%%%%

\renewcommand{\familydefault}{\rmdefault}


\newcommand{\mycolor}[1] {
\setbeamercolor{normal text}{fg={#1}} % Recolors standard body text
\setbeamercolor{structure}{fg={#1}}   % Recolors bullets and structure items
\usebeamercolor[fg]{normal text}      % Applies the normal text color immediately
}


\usepackage{amsmath}
\usepackage{mathtools}

\usepackage{cancel}
\usepackage{centernot}

\usepackage{changepage}
\usepackage{subcaption}

\usepackage{qrcode}

\newcommand{\bo}[1] {\mathbf{#1}}
\newcommand{\R}{\mathbb{R}} 
\newcommand{\T}{^\top}     



\newcommand{\mydate}{Fall 2026}

\newcommand{\mygit}{\textcolor{blue}{\href{https://github.com/SergeiSa/Advanced-Control-2026}{github.com/SergeiSa/Advanced-Control-2026}}}

\newcommand{\myqr}{ \textcolor{black}{\qrcode[height=1.5in]{https://github.com/SergeiSa/Advanced-Control-2026}}
}




\newcommand{\bref}[2] {\textcolor{blue}{\href{#1}{#2}}}

%%%%%%%%%%%%%%%%%%%%%%%%%%%%
% code settings

\usepackage{listings}
\usepackage{color}
% \definecolor{mygreen}{rgb}{0,0.6,0}
% \definecolor{mygray}{rgb}{0.5,0.5,0.5}
\definecolor{mymauve}{rgb}{0.58,0,0.82}
\lstset{ 
	backgroundcolor=\color{white},   % choose the background color; you must add \usepackage{color} or \usepackage{xcolor}; should come as last argument
	basicstyle=\footnotesize,        % the size of the fonts that are used for the code
	breakatwhitespace=false,         % sets if automatic breaks should only happen at whitespace
	breaklines=true,                 % sets automatic line breaking
	captionpos=b,                    % sets the caption-position to bottom
	commentstyle=\color{mygreen},    % comment style
	deletekeywords={...},            % if you want to delete keywords from the given language
	escapeinside={\%*}{*)},          % if you want to add LaTeX within your code
	extendedchars=true,              % lets you use non-ASCII characters; for 8-bits encodings only, does not work with UTF-8
	firstnumber=0000,                % start line enumeration with line 0000
	frame=single,	                   % adds a frame around the code
	keepspaces=true,                 % keeps spaces in text, useful for keeping indentation of code (possibly needs columns=flexible)
	keywordstyle=\color{blue},       % keyword style
	language=Octave,                 % the language of the code
	morekeywords={*,...},            % if you want to add more keywords to the set
	numbers=left,                    % where to put the line-numbers; possible values are (none, left, right)
	numbersep=5pt,                   % how far the line-numbers are from the code
	numberstyle=\tiny\color{mygray}, % the style that is used for the line-numbers
	rulecolor=\color{black},         % if not set, the frame-color may be changed on line-breaks within not-black text (e.g. comments (green here))
	showspaces=false,                % show spaces everywhere adding particular underscores; it overrides 'showstringspaces'
	showstringspaces=false,          % underline spaces within strings only
	showtabs=false,                  % show tabs within strings adding particular underscores
	stepnumber=2,                    % the step between two line-numbers. If it's 1, each line will be numbered
	stringstyle=\color{mymauve},     % string literal style
	tabsize=2,	                   % sets default tabsize to 2 spaces
	title=\lstname                   % show the filename of files included with \lstinputlisting; also try caption instead of title
}


%%%%%%%%%%%%%%%%%%%%%%%%%%%%
% URL settings
\hypersetup{
	colorlinks=false,
	linkcolor=blue,
	filecolor=blue,      
	urlcolor=blue,
}

%%%%%%%%%%%%%%%%%%%%%%%%%%

%%%%%%%%%%%%%%%%%%%%%%%%%%%%
% tikz settings

\usepackage{tikz}
\tikzset{every picture/.style={line width=0.75pt}}
\usetikzlibrary{shapes.geometric, decorations.markings, positioning, patterns, calc, arrows.meta}

\newcommand{\myqrframe}{
	\begin{frame}
		\begin{center}
			\centerline{Слайды доступны на Github:}
			\bigskip
			\centerline{\mygit}
			\bigskip
			\myqr
		\end{center}
	\end{frame}
}


\subtitle{Современные методы управления роботами}
\author{Сергей Савин}
\date{2026}


\title{Стохастическое уравнение Беллмана}
\begin{document}
	\maketitle
	
	
	\begin{frame}{Содержание}
		\begin{itemize}
			\item Стохастические последовательности и математическое ожидание
			\item Стохастические линейные системы
			\item Конечногoризонтная стохастическая стоимость
			\item Стохастическое уравнение Беллмана
			\item Стохастическое уравнение Риккати
			\item Бесконечногoризонтная стоимость с дисконтированием
			\item Функции стоимости и функции ценности действия
			\item Уравнения Беллмана и Риккати с дисконтированием
			\item Приложение: вывод стохастического уравнения Риккати
		\end{itemize}
	\end{frame}
	
	
	
	%------------------------------------------------
	\begin{frame}{Стохастические последовательности}
		\begin{flushleft}
			
			{\small 
			Пусть $w_k$ — случайная (стохастическая) последовательность с математическим ожиданием
			\[
			\mu = \mathbb{E}[w_k].
			\]
			
			По $N$ выборкам можно оценить её математическое ожидание с помощью выборочного среднего:
			\[
			\hat{\mu}_N = \frac{1}{N}\sum_{i=k}^{k+N-1} w_i.
			\]
			
			При соответствующих предположениях (например, независимые одинаково распределённые
			выборки с конечным средним) имеем
			\[
			\lim_{N\to\infty}\hat{\mu}_N = \mathbb{E}[w_k].
			\]
			
			
			Математическое ожидание — линейный оператор. В частности, для постоянных матриц
			$A,B$ и случайных векторов $x,u$:
			\[
			\boxed{
				\mathbb{E}[Ax+Bu]
				= A\mathbb{E}[x]+B\mathbb{E}[u]
			}
			\]
		}
		\end{flushleft}
	\end{frame}
	%------------------------------------------------
	
	\begin{frame}{Стохастические линейные системы}
		\begin{flushleft}
			
			Рассмотрим дискретную LTI-систему:
			\[
			x_{k+1}=Ax_k+w_k,
			\]
			где $w_k$ — элемент стохастической последовательности.
			
			Можно сделать следующие наблюдения:
			
			\begin{itemize}
				\item Не зная реализации $w_k$, мы не можем предсказать
				точное значение $x_{k+1}$.
				
				\item Следовательно, даже при известном начальном состоянии $x_{k+1}$
				является случайной величиной.
				
				\item Хотя точные значения состояния предсказать нельзя, можно
				анализировать их статистические свойства.
			\end{itemize}
			
			\vspace{0.2cm}
			
			Например, используя линейность математического ожидания:
			\[
			\boxed{
				\mathbb{E}[x_{k+1}]
				= A\mathbb{E}[x_k]+\mathbb{E}[w_k]
			}
			\]
			
		\end{flushleft}
	\end{frame}
	%------------------------------------------------
	
	\begin{frame}{Конечногoризонтная стохастическая стоимость}
		\begin{flushleft}
			
			Введём квадратичную функцию стоимости, подходящую для
			стохастической линейной системы:
			\[
			J = \mathbb{E}\left[
			\sum_{k=0}^{N-1} \ell(x_k,u_k)
			+ x_N^\top Q_N x_N
			\right],
			\]
			где стоимость этапа
			\[
			\ell(x_k,u_k)
			= x_k^\top Q x_k + u_k^\top R u_k.
			\]
			
			Здесь $Q\succeq 0, Q_N\succeq 0$ и $R\succ 0$.
			
			\vspace{0.3cm}
			
			Как и в детерминированном случае, определим функцию стоимости
			на шаге $k$ как минимальную ожидаемую стоимость до конца:
			\[
			\boxed{
				V_k(x) =
				\min_{u_k,\ldots,u_{N-1}}
				\mathbb{E}\left[
				\sum_{i=k}^{N-1}\ell(x_i,u_i)
				+x_N^\top Q_N x_N
				\,\middle|\, x_k=x
				\right]
			}
			\]
			
		\end{flushleft}
	\end{frame}
	%------------------------------------------------
	
	
	
	
	%------------------------------------------------
	\begin{frame}{Стохастическое уравнение Беллмана}
		\begin{flushleft}
			
			Как и в детерминированном случае, можно использовать принцип Беллмана,
			чтобы выразить функцию стоимости рекуррентно:
			
			\[
			V_k(x_k) =
			\min_{u_k}
			\left\{
			\ell(x_k,u_k)
			+ \mathbb{E}\left[
			V_{k+1}(x_{k+1}) \mid x_k,u_k
			\right]
			\right\}.
			\]
			
			Для нашей стохастической линейной системы:
			\[
			x_{k+1}=Ax_k+Bu_k+w_k,
			\]
			получаем
			\[
			\boxed{
				V_k(x) =
				\min_u
				\left\{
				x^\top Qx+u^\top Ru
				+\mathbb{E}\left[
				V_{k+1}(Ax+Bu+w_k)
				\right]
				\right\}.
			}
			\]
			
			Терминальное условие:
			\[
			V_N(x)=x^\top Q_Nx.
			\]
			
		\end{flushleft}
	\end{frame}
	%------------------------------------------------
	
	\begin{frame}{Стохастическое уравнение Риккати}
		\begin{flushleft}
			Для невозмущённых независимых возмущений с нулевым средним и ковариацией
			\[
			\mathbb{E}[w_k]=0,\qquad
			\mathbb{E}[w_kw_k^\top]=W,
			\]
			функция стоимости сохраняет квадратичную форму:
			\[
			\boxed{
				V_k(x)=x^\top P_kx+c_k.
			}
			\]
			
			Подстановка этого выражения в уравнение Беллмана даёт
			оптимальный линейный закон управления:
			\[
			u_k=K_kx_k,
			\]
			где
			\[
			K_k =
			-(R+B^\top P_{k+1}B)^{-1}
			B^\top P_{k+1}A.
			\]
			
			Матрица $P_k$ находится с помощью рекурсии Риккати:
			\[
			P_k=Q+A^\top P_{k+1}A
			-A^\top P_{k+1}B
			(R+B^\top P_{k+1}B)^{-1}
			B^\top P_{k+1}A.
			\]
			
			с терминальным условием $P_N=Q_N$.
		\end{flushleft}
	\end{frame}
	%------------------------------------------------
	
	
	\begin{frame}{Бесконечногoризонтная стоимость}
		\begin{flushleft}
			
				Можно рассмотреть бесконечногoризонтную версию стохастической
				задачи оптимального управления.
				
				\bigskip
				
				Одна из проблем в том, что накопленная стоимость может стать бесконечной:
				даже если система устойчива, постоянные стохастические возмущения
				продолжают вносить неопределённость в состояние.
				
				
		\end{flushleft}
	\end{frame}
	
	
	
	\begin{frame}{Бесконечногoризонтная стоимость}
		\begin{flushleft}
			
				
				Один из способов справиться с этой проблемой — ввести коэффициент
				дисконтирования $0 < \gamma < 1$:
				
				\[
				\boxed{
					J = \mathbb{E}\left[
					\sum_{k=0}^{\infty}
					\gamma^k \ell(x_k,u_k)
					\right].
				}
				\]
				
				\vspace{0.2cm}
				
				\begin{itemize}
					\item Дисконтирование позволяет ожидаемой накопленной стоимости оставаться
					конечной при стандартных предположениях об ограниченности стоимости.
					
					\item Оно отдаёт приоритет немедленным затратам, постепенно
					уменьшая важность затрат в более отдалённом будущем.
					
					\item Как и прежде, бесконечногoризонтная постановка не содержит терминальной
					составляющей стоимости.
				\end{itemize}
			
		\end{flushleft}
	\end{frame}
	
	
	
	\begin{frame}{Функция стоимости с дисконтированной стоимостью}
		\begin{flushleft}
			
			Для заданной политики $\pi$ определим функцию стоимости состояния
			на политике как ожидаемую дисконтированную стоимость до конца:
			
			\[
			V^\pi(x) =
			\mathbb{E}_\pi\left[
			\sum_{k=0}^{\infty}
			\gamma^k \ell(x_k,u_k)
			\,\middle|\, x_0=x
			\right].
			\]
			
			\vspace{0.3cm}
			
			Оптимальная функция стоимости состояния — минимум по всем
			допустимым политикам:
			
			\[
			V^*(x)=\min_\pi V^\pi(x).
			\]
			
		\end{flushleft}
	\end{frame}
	
	\begin{frame}{Уравнение Беллмана с дисконтированной стоимостью}
		\begin{flushleft}
			
			Уравнение оптимальности Беллмана для бесконечного горизонта имеет вид:
			%
			\[
			V^*(x)=
			\min_u\left\{
			\ell(x,u)+
			\gamma\mathbb{E}\left[
			V^*(x_{k+1})\mid x_k=x,u_k=u
			\right]
			\right\}.
			\]
			%
			Можно ввести функцию ценности действия $\mathcal{Q}^\pi(x,u)$,
			которая описывает ожидаемую стоимость до конца, если применить действие $u$
			в состоянии $x$, а затем следовать политике $\pi$:
			%
			\[
			\mathcal{Q}^\pi(x,u)=
			\mathbb{E}_\pi\left[
			\sum_{k=0}^{\infty}\gamma^k\ell(x_k,u_k)
			\,\middle|\,
			x_0=x,\ u_0=u
			\right].
			\]
			%
			Функции стоимости состояния и ценности действия связаны соотношением:
			%
			\[
			V^\pi(x)=
			\mathbb{E}_{u\sim\pi(\cdot|x)}
			[\mathcal{Q}^\pi(x,u)].
			\]
			%
			Для оптимальной политики: $V^*(x)=\min_u \mathcal{Q}^*(x,u).$
			%
		\end{flushleft}
	\end{frame}
	
	
	\begin{frame}{Уравнение Риккати с дисконтированной стоимостью}
		\begin{flushleft}
			
			Для стохастического LQR оптимальная функция стоимости по-прежнему квадратична:
			\[
			V^*(x)=x^\top Px+c.
			\]
			
			Подстановка этого выражения в уравнение Беллмана даёт
			дисконтированное алгебраическое уравнение Риккати:
			
			\[
			P=Q+\gamma A^\top PA
			-\gamma^2 A^\top PB
			(R+\gamma B^\top PB)^{-1}
			B^\top PA.
			\]
			
			Оптимальный закон управления линеен:
			\[
			u^*=-Kx,\qquad
			K=(R+\gamma B^\top PB)^{-1}
			\gamma B^\top PA.
			\]
			
			Постоянный член учитывает стохастические возмущения:
			\[
			c=\frac{\gamma}{1-\gamma}
			\operatorname{tr}(PW).
			\]
			
		\end{flushleft}
	\end{frame}
	
	
	
	\begin{frame}{Литература}
		\begin{itemize}
			
			\item MIT OpenCourseWare:
			\bref{https://ocw.mit.edu/courses/16-323-principles-of-optimal-control-spring-2008/pages/lecture-notes/}
			{Principles of Optimal Control: Lecture 3, Dynamic Programming and Discrete LQR.}
			
			\item Russ Tedrake -- MIT:
			\bref{https://underactuated.mit.edu/lqr.html}
			{Underactuated Robotics: Linear Quadratic Regulators, including the discrete-time Riccati equation.}
			
			\item University of Washington:
			\bref{https://faculty.washington.edu/ratliffl/teaching/2020_LinearSystems.pdf}
			{Lecture 10: Linear Quadratic Regulator -- Discrete Time.}
			
			\item SciPy documentation:
			\bref{https://docs.scipy.org/doc/scipy/reference/generated/scipy.linalg.solve_discrete_are.html}
			{Linear Algebra: \texttt{solve\_discrete\_are}, numerical solution of the discrete-time algebraic Riccati equation.}
			
		\end{itemize}
	\end{frame}
	
	
	\myqrframe
	
	\begin{frame}{Приложение: вывод стохастического уравнения Риккати}
		\begin{flushleft}
			Подставим квадратичную функцию стоимости в уравнение Беллмана:
			\[
			\begin{aligned}
				V_k(x)=\min_u\big\{&
				x^\top Qx+u^\top Ru\\
				&+\mathbb{E}[
				(Ax+Bu+w_k)^\top P_{k+1}
				(Ax+Bu+w_k)]\\
				&+c_{k+1}\big\}.
			\end{aligned}
			\]
			
			Используя $\mathbb{E}[w_k]=0$ и
			$\mathbb{E}[w_kw_k^\top]=W$:
			\[
			\begin{aligned}
				V_k(x)=\min_u\big\{&
				x^\top Qx+u^\top Ru\\
				&+(Ax+Bu)^\top P_{k+1}(Ax+Bu)\\
				&+\operatorname{tr}(P_{k+1}W)+c_{k+1}
				\big\}.
			\end{aligned}
			\]
			
			
			Приравнивая производную по $u$ к нулю, получаем
			\[
			u_k=-(R+B^\top P_{k+1}B)^{-1}
			B^\top P_{k+1}Ax_k.
			\]
			
		\end{flushleft}
	\end{frame}
	
	\begin{frame}{Приложение: рекурсия Риккати}
		\begin{flushleft}
			
			Подставляя оптимальный закон управления обратно в уравнение Беллмана,
			получаем:
			\[
			\begin{aligned}
				P_k={}&Q+A^\top P_{k+1}A\\
				&-A^\top P_{k+1}B
				(R+B^\top P_{k+1}B)^{-1}
				B^\top P_{k+1}A,
			\end{aligned}
			\]
			
			и
			\[
			c_k=c_{k+1}+\operatorname{tr}(P_{k+1}W).
			\]
			
			Терминальные условия:
			\[
			P_N=Q_N,\qquad c_N=0.
			\]
			
		\end{flushleft}
	\end{frame}
	
	
	
	
	
\end{document}